\documentclass[11pt,a4paper,fontset=fandol]{ctexart}

\usepackage[a4paper,top=2.35cm,bottom=2.55cm,left=2.3cm,right=2.3cm,headheight=18pt,headsep=0.55cm,footskip=26pt]{geometry}
\usepackage{amsmath,amssymb,amsthm}
\usepackage{fancyhdr}
\usepackage{lastpage}
\usepackage{enumitem}
\usepackage{titlesec}
\usepackage{xcolor}
\usepackage{tikz}
\usepackage{hyperref}
\usepackage{bookmark}
\usepackage[most]{tcolorbox}
\usepackage{setspace}
\usepackage{array}
\usetikzlibrary{arrows.meta,positioning}

\hypersetup{
  colorlinks=true,
  linkcolor=blue!50!black,
  urlcolor=blue!50!black,
  citecolor=blue!50!black
}

\setstretch{1.15}
\setlength{\parindent}{2em}
\setlength{\parskip}{0.28em}
\allowdisplaybreaks
\setlist[enumerate]{itemsep=0.35em, topsep=0.35em, leftmargin=2.2em}
\setlist[itemize]{itemsep=0.25em, topsep=0.25em, leftmargin=1.9em}

\definecolor{mainblue}{RGB}{36,83,140}
\definecolor{lightblue}{RGB}{241,246,252}
\definecolor{lightgray}{RGB}{246,246,246}
\definecolor{accent}{RGB}{153,76,0}
\definecolor{rulegray}{RGB}{110,118,128}

\titleformat{\section}
  {\Large\bfseries\color{mainblue}}
  {\thesection.}
  {0.45em}
  {}
  [\vspace{0.25em}\color{rulegray}\titlerule]
\titlespacing*{\section}{0pt}{1.2em}{0.8em}
\titleformat{\subsection}{\large\bfseries\color{mainblue}}{\thesubsection}{0.5em}{}

\pagestyle{fancy}
\fancyhf{}
\fancyhead[L]{\small 函数图像平移提高训练题单}
\fancyhead[C]{\small 代数专题：图像平移与方程变形}
\fancyhead[R]{\small 刘文博}
\fancyfoot[C]{\small 第 \thepage 页 / 共 \pageref{LastPage} 页}
\renewcommand{\headrulewidth}{0.55pt}
\renewcommand{\footrulewidth}{0.35pt}

\newtcolorbox{goalbox}[1]{
  breakable,
  colback=lightblue,
  colframe=mainblue,
  boxrule=0.7pt,
  arc=1mm,
  title=\textbf{#1},
  fonttitle=\bfseries
}

\newtcolorbox{notebox}[1]{
  breakable,
  colback=lightgray,
  colframe=black!50,
  boxrule=0.55pt,
  arc=1mm,
  title=\textbf{#1},
  fonttitle=\bfseries
}

\newtheoremstyle{formalexample}
  {0.8em}{0.8em}{\normalfont}{}{\bfseries\color{accent}}{．}{0.6em}{}
\theoremstyle{formalexample}
\newtheorem{exercise}{题目}[section]
\newenvironment{solution}{\par\noindent\textbf{解：}}{\hfill$\square$\par}

\begin{document}

\thispagestyle{empty}
\begin{titlepage}
\centering
\vspace*{0.9cm}
\begin{tcolorbox}[
  enhanced,
  width=0.92\textwidth,
  colback=white,
  colframe=mainblue,
  boxrule=1pt,
  arc=0mm,
  left=6mm,
  right=6mm,
  top=8mm,
  bottom=8mm
]
\centering
{\fontsize{24}{30}\selectfont\bfseries 函数图像平移提高训练题单\par}
\vspace{0.6cm}
{\fontsize{17}{22}\selectfont 适合初中阶段函数入门与巩固训练\par}

\vspace{1cm}
\begin{tcolorbox}[colback=lightblue,colframe=mainblue,width=0.84\textwidth,boxrule=1pt]
\centering
{\large 训练目标：把“会背公式”提升到“会判断、会变形、会检验”}\\[0.35em]
{\normalsize 建议分 2--3 次完成，先独立做题，再对照详细解答订正}
\end{tcolorbox}

\vspace{1cm}
\renewcommand{\arraystretch}{1.42}
\begin{tabular}{>{\bfseries}p{3.5cm}p{8cm}}
题目数量： & 16 题，按层次递进 \\
专题重点： & 向上、向下、向左、向右平移；综合平移；反向判断；关键点检验 \\
答案形式： & 末尾附详细解答，强调“为什么这样改方程” \\
编译方式： & XeLaTeX \\
\end{tabular}

\vspace{0.9cm}
{\large \today\par}
\end{tcolorbox}
\end{titlepage}

\tableofcontents
\newpage

\section{使用建议}

\begin{goalbox}{做函数平移题时，优先问自己的 4 个问题}
\begin{itemize}
  \item 这是上下平移，还是左右平移？
  \item 如果是上下平移，我是不是应该在函数\textbf{外面}加减？
  \item 如果是左右平移，我是不是应该在函数\textbf{里面}改 $x$？
  \item 我能不能用顶点、截距点或特殊点检验结果？
\end{itemize}
\end{goalbox}

\begin{notebox}{建议计时}
\begin{itemize}
  \item 基础题：每题 2--4 分钟；
  \item 综合题：每题 5--8 分钟；
  \item 反向判断题：先写成 $y=f(x-a)+b$ 的统一形式再判断。
\end{itemize}
\end{notebox}

\section{基础巩固}

\begin{center}
\begin{tikzpicture}[scale=0.72,>=Stealth]
  \draw[->,gray!70] (-4.2,0) -- (4.2,0) node[right] {$x$};
  \draw[->,gray!70] (0,-2.2) -- (0,5.8) node[above] {$y$};
  \draw[smooth,domain=-2.1:2.1,samples=100,thick,mainblue] plot (\x,{abs(\x)});
  \draw[smooth,domain=-4.1:0.1,samples=100,thick,red!70!black] plot (\x,{abs(\x+2)-1});
  \fill[mainblue] (0,0) circle (2pt);
  \fill[red!70!black] (-2,-1) circle (2pt);
  \node[mainblue] at (2.4,2.4) {$y=|x|$};
  \node[red!70!black] at (-2.6,2.2) {$y=|x+2|-1$};
  \node[draw=none,fill=none,font=\small\color{rulegray}] at (0,-1.7) {关键点 $(0,0)$ 平移到 $(-2,-1)$，结果就更容易判断};
\end{tikzpicture}
\end{center}

\begin{exercise}
把函数
\[
y=x^2
\]
向上平移 $5$ 个单位，写出新方程。
\end{exercise}

\begin{exercise}
把函数
\[
y=|x|
\]
向左平移 $2$ 个单位，再向下平移 $1$ 个单位，写出新方程。
\end{exercise}

\begin{exercise}
把函数
\[
y=2x-5
\]
向右平移 $4$ 个单位，写出新方程。
\end{exercise}

\begin{exercise}
把函数
\[
y=\frac{1}{x}
\]
向右平移 $2$ 个单位，再向上平移 $3$ 个单位，写出新方程。
\end{exercise}

\section{反向判断}

\begin{exercise}
已知函数
\[
y=(x-3)^2+2,
\]
说出它是由
\[
y=x^2
\]
怎样平移得到的。
\end{exercise}

\begin{exercise}
已知函数
\[
y=|x+4|-3,
\]
说出它是由
\[
y=|x|
\]
怎样平移得到的。
\end{exercise}

\begin{exercise}
已知函数
\[
y=2(x-1)+5,
\]
说出它是由
\[
y=2x
\]
怎样平移得到的。
\end{exercise}

\begin{exercise}
已知函数
\[
y=f(x+2)-7,
\]
说出它是由
\[
y=f(x)
\]
怎样平移得到的。
\end{exercise}

\section{综合应用}

\begin{exercise}
把函数
\[
y=x^2-2x
\]
向右平移 $3$ 个单位，再向上平移 $1$ 个单位，写出化简后的新方程。
\end{exercise}

\begin{exercise}
把函数
\[
y=2x+3
\]
向左平移 $2$ 个单位，再向下平移 $6$ 个单位，写出化简后的新方程。
\end{exercise}

\begin{exercise}
已知函数
\[
y=|x|
\]
经过平移后，顶点变成 $(-3,2)$。求新方程。
\end{exercise}

\begin{exercise}
把函数
\[
y=(x-4)^2+5
\]
向左平移 $1$ 个单位，再向下平移 $7$ 个单位，写出新方程。
\end{exercise}

\begin{exercise}
原函数
\[
y=x^2
\]
上的点 $A(1,1)$，在图像向右平移 $2$ 个单位、向下平移 $3$ 个单位后，对应点 $A'$ 的坐标是什么？
\end{exercise}

\begin{exercise}
已知抛物线
\[
y=(x-1)^2
\]
经过平移后，它的顶点从 $(1,0)$ 移到 $(4,-2)$。求新方程。
\end{exercise}

\begin{exercise}
已知原图像上点 $(2,5)$ 在函数
\[
y=f(x)
\]
上。平移后该点对应到新图像上的点 $(-1,9)$。求新图像的方程。
\end{exercise}

\begin{exercise}
已知函数
\[
y=x^2+6x+11,
\]
说出它是由
\[
y=x^2
\]
怎样平移得到的。
\end{exercise}

\newpage
\section{常见易错点总结}

\begin{goalbox}{做函数平移题时，最容易丢分的 5 个地方}
\begin{enumerate}
  \item \textbf{把左右平移的符号写反。}  
  向左平移时应写成 $x+a$，向右平移时应写成 $x-a$。

  \item \textbf{只改了一部分，没有改整个 $x$。}  
  例如 $y=x^2-2x$ 向右平移后，必须把整个式子里的 $x$ 都换成 $x-3$。

  \item \textbf{看到平移就立刻化简，结果中途丢项。}  
  建议先保留括号，再慢慢展开。

  \item \textbf{只会正向写，不会反向看。}  
  反向题要先整理成 $y=f(x-a)+b$ 的形式。

  \item \textbf{不拿关键点检验。}  
  例如顶点、截距点、特殊点跟着一起移动，最能帮助你发现符号错误。
\end{enumerate}
\end{goalbox}

\begin{notebox}{一个实用检查顺序}
\begin{itemize}
  \item 先分清上下还是左右；
  \item 再决定是改“外面”还是改“里面”；
  \item 最后用关键点检查结果是否合理。
\end{itemize}
\end{notebox}

\newpage
\appendix
\section{详细解答}

\subsection*{基础巩固}

\textbf{1}\begin{solution}
向上平移 $5$ 个单位，就是在原函数外面加 $5$，所以新方程为
\[
\boxed{y=x^2+5}.
\]
\end{solution}

\textbf{2}\begin{solution}
先向左平移 $2$ 个单位，把 $x$ 换成 $x+2$，得到
\[
y=|x+2|.
\]
再向下平移 $1$ 个单位，得到
\[
\boxed{y=|x+2|-1}.
\]
\end{solution}

\textbf{3}\begin{solution}
向右平移 $4$ 个单位，应把 $x$ 换成 $x-4$，所以
\[
y=2(x-4)-5.
\]
化简得
\[
\boxed{y=2x-13}.
\]
\end{solution}

\textbf{4}\begin{solution}
先向右平移 $2$ 个单位，得到
\[
y=\frac{1}{x-2}.
\]
再向上平移 $3$ 个单位，得到
\[
\boxed{y=\frac{1}{x-2}+3}.
\]
\end{solution}

\subsection*{反向判断}

\textbf{5}\begin{solution}
方程
\[
y=(x-3)^2+2
\]
已经写成了“右 $3$、上 $2$”的形式，所以它是由 $y=x^2$ \textbf{向右平移 $3$ 个单位，再向上平移 $2$ 个单位}得到的。
\end{solution}

\textbf{6}\begin{solution}
在
\[
y=|x+4|-3
\]
中，式子里的 $x+4$ 表示向左平移 $4$ 个单位，外面的 $-3$ 表示向下平移 $3$ 个单位。

所以它是由 $y=|x|$ \textbf{向左平移 $4$ 个单位，再向下平移 $3$ 个单位}得到的。
\end{solution}

\textbf{7}\begin{solution}
把
\[
y=2(x-1)+5
\]
和统一形式
\[
y=f(x-a)+b
\]
比较，可知 $a=1$，$b=5$。

所以它是由 $y=2x$ \textbf{向右平移 $1$ 个单位，再向上平移 $5$ 个单位}得到的。
\end{solution}

\textbf{8}\begin{solution}
方程
\[
y=f(x+2)-7
\]
可看成
\[
y=f(x-(-2))+(-7).
\]
因此图像是由 $y=f(x)$ \textbf{向左平移 $2$ 个单位，再向下平移 $7$ 个单位}得到的。
\end{solution}

\subsection*{综合应用}

\textbf{9}\begin{solution}
向右平移 $3$ 个单位，先把整个式子里的 $x$ 换成 $x-3$：
\[
y=(x-3)^2-2(x-3).
\]
再向上平移 $1$ 个单位：
\[
y=(x-3)^2-2(x-3)+1.
\]
展开化简：
\[
y=x^2-6x+9-2x+6+1=x^2-8x+16.
\]
所以新方程为
\[
\boxed{y=x^2-8x+16}.
\]
\end{solution}

\textbf{10}\begin{solution}
先向左平移 $2$ 个单位：
\[
y=2(x+2)+3.
\]
再向下平移 $6$ 个单位：
\[
y=2(x+2)+3-6.
\]
化简得
\[
y=2x+1.
\]
所以新方程为
\[
\boxed{y=2x+1}.
\]
\end{solution}

\textbf{11}\begin{solution}
函数 $y=|x|$ 的顶点是 $(0,0)$。

平移后顶点变成 $(-3,2)$，说明图像向左平移了 $3$ 个单位，再向上平移了 $2$ 个单位。

因此新方程为
\[
\boxed{y=|x+3|+2}.
\]
\end{solution}

\textbf{12}\begin{solution}
先向左平移 $1$ 个单位，应把原式中的 $x$ 换成 $x+1$：
\[
y=((x+1)-4)^2+5=(x-3)^2+5.
\]
再向下平移 $7$ 个单位：
\[
y=(x-3)^2+5-7=(x-3)^2-2.
\]
所以新方程为
\[
\boxed{y=(x-3)^2-2}.
\]
\end{solution}

\textbf{13}\begin{solution}
点 $A(1,1)$ 跟着图像一起平移。

向右平移 $2$ 个单位后，横坐标加 $2$，变成 $3$；

向下平移 $3$ 个单位后，纵坐标减 $3$，变成 $-2$。

所以
\[
\boxed{A'(3,-2)}.
\]
\end{solution}

\textbf{14}\begin{solution}
原抛物线 $y=(x-1)^2$ 的顶点是 $(1,0)$。

现在顶点移到 $(4,-2)$，说明图像向右平移了 $3$ 个单位，再向下平移了 $2$ 个单位。

因此新方程为
\[
\boxed{y=(x-4)^2-2}.
\]
\end{solution}

\textbf{15}\begin{solution}
原图像上点 $(2,5)$ 平移后变成 $(-1,9)$。

横坐标从 $2$ 变到 $-1$，减少了 $3$，说明向左平移 $3$ 个单位；

纵坐标从 $5$ 变到 $9$，增加了 $4$，说明向上平移 $4$ 个单位。

因此新图像方程为
\[
\boxed{y=f(x+3)+4}.
\]
\end{solution}

\textbf{16}\begin{solution}
先配方：
\[
y=x^2+6x+11=(x+3)^2+2.
\]
因此它是由 $y=x^2$ \textbf{向左平移 $3$ 个单位，再向上平移 $2$ 个单位}得到的。
\end{solution}

\end{document}
